\section{Closed routes, with the mechanism that closed them}

This section is a graveyard with load-bearing headstones. Each entry below is a route that was
tried against the $\sqrt{}$-escape wall (Part~III) and failed, but the way it failed is itself a
theorem, a computation, or a structural identity that survives the failure. A closed route is not
a blank; it is a constraint on every future attempt. We record, for each closure: the route as
conceived, the exact mechanism that closed it, the precise scope of the closure (what it rules
out and what it leaves untouched), and whatever machinery salvages out of the wreckage. We also
flag, for every closure, whether the exclusion is a statement about the \emph{object} (the real
number, the achievement set, the actual orbit) or about a \emph{representation} of it (a
coordinate, a finite-window signature, a proof-route's hypothesis set). Confusing the two is the
single most common way a closed route gets silently re-opened by a later attempt using different
words for the same representation.

Throughout, $\mathrm{rem}(s)$ is the seam-row remainder, $\lambda_n := \Delta_n/2^n$ the normalized seam
deviation (Part~I, SE-3), and $\mathcal A$ the Mersenne achievement set. Reset $\sqrt{}$-escape
from row 10 is the statement $|\mathrm{rem}(r+1)-2^{r+1}| > 2^{(r+5)/2}$ for every reset row $r\ge 10$
(RC-4, Part~III); this is the wall every closure below is measured against.

\subsection{Four closed attack families on the wall itself}

\begin{obs}[Prime-doubling — closed by explicit counterexample]
\label{obs:prime-doubling}
\textbf{Route as conceived.} The universal recursion $K(M+1) = 2K(M) - (\tau(M+1)-1)$ (Part~I,
SE-1) doubles at every step and subtracts the divisor-excess $\tau(M+1)-1$. At a prime $M+1=p$,
$\tau(p)=2$ so the subtracted term is $1$; the naive hope is that primality of an index forces the
step to be "almost pure doubling," and that prime density (Chebyshev/PNT-scale) then bounds how
long a run of near-pure-doubling steps can last, which would bound the danger run-length $L_r$ in
SE-3's equivalence (III).

\textbf{Exact mechanism that closed it.} The hope requires primality to control the recursion
\emph{cumulatively}, but composite indices subtract more than primes add back, so a long run can
mix primes and composites freely without breaking. Certified witness: the run at $n=607$ spans
$M\in(607,617]$ with the drift quantity $J$ staying in $[1,5]$ for eleven consecutive steps, and
three of the indices inside that window — $607,613,617$ — are themselves prime. More broadly, of
the 40 longest runs found for $n\le 1500$, 26 contain a prime index somewhere inside them. Primality
of an index is simply uncorrelated with the recursion staying quiet.

\textbf{Scope.} Closes: any argument that tries to derive a run-length bound on SE-3's $L_r$, or a
bound on the drift $J$ in the universal recursion, from prime-counting alone (density,
Chebyshev bounds, PNT-scale gap statistics). Does not touch: arguments using the actual value of
$\tau(M+1)-1$ (not just whether $M+1$ is prime), or any argument working at the divisor-sum level
rather than the prime/composite dichotomy.
\coord{prime-density}

\textbf{Object or representation.} About the \emph{object}: the counterexample is a concrete
witness inside the actual $K(M)$ recursion, not an artifact of how the recursion is coordinatized.

\textbf{Salvage.} None as a proof tool — but the witness itself (run length 11 at $n=607$,
26/40 longest runs prime-touched) is now a banked disconfirming data point: any future
prime-density-flavoured attack on run-length should be checked against this witness first.
\ev{Cert}
\end{obs}

\begin{obs}[Parity/perfect-squares — closed as a tautology of the recursion shape]
\label{obs:parity-squares}
\textbf{Route as conceived.} $K(M)$ is even exactly when $M$ is a perfect square (a true,
provable fact). The hope was to use this parity law to forbid certain residues or run patterns in
the digit stream of $K$, since parity constraints are a classical first attack on digit-run
questions.

\textbf{Exact mechanism that closed it.} The recursion $J(M+1) = 2J(M) - \delta(M+1)$ makes
$J(M+1)\equiv\delta(M+1)\pmod 2$ an \emph{identity for any digit stream whatsoever} that this
recursion shape produces — it holds regardless of what $\delta$ counts, Mersenne-specific or not.
The "$K(M)$ even iff $M$ a perfect square" fact is real, but it is a restatement of the recursion's
algebraic shape, not new information: it forbids no residue and no run length. It is vacuous as an
obstruction because it would hold for \emph{every} recursion of the form $J(M+1)=2J(M)-\delta(M+1)$,
Mersenne or not.

\textbf{Scope.} Closes: the parity/perfect-square route completely and permanently, for this
recursion shape — not a partial closure. Also closes it in advance for any structurally identical
recursion (e.g.\ a hypothetical \#249 digit-carry recursion of the same doubling-minus-forcing
shape), since the vacuity is a fact about the shape, not about $\tau$ or Mersenne weights.
\coord{parity-tautology}

\textbf{Object or representation.} About a \emph{representation}: the identity is a fact about the
recursion's algebraic form (base-2 doubling-with-forcing), not about the value $K(M)$ or the
target constant $C=E-3/2$ itself. It is empty precisely because it never touches the object.

\textbf{Salvage.} The identity remains true and is a useful sanity check (any proposed digit-run
claim about $K$ must be consistent with it), but it supplies zero constraining power on its own.
\ev{Math}
\end{obs}

\begin{obs}[Irrationality-measure — structurally excluded by an exponent budget]
\label{obs:irrationality-measure}
\textbf{Route as conceived.} The wall reduces (SE-2) to: the greedy skip-set sum
$\sum_{d\in\mathrm{Skip}_n}x_d$ never approximates $C:=E-\tfrac32=0.1066951524152917\ldots$ (the shifted
Erdős–Borwein constant) to within $\sim 2^{-1.5n}$. Since $C$ has a known irrationality-measure
upper bound, the natural hope is that a Diophantine-approximation exponent for $C$ directly
supplies the needed non-approximability at scale $n$.

\textbf{Exact mechanism that closed it.} The skip-set $\mathrm{Skip}_n$ has size $|\mathrm{Skip}_n|\approx n^2/4$ —
quadratic in $n$, not linear. An irrationality-measure exponent $\mu$ for $C$ bounds how well a
\emph{single} rational with denominator $q$ can approximate $C$; specializing to a rational built
from $\approx n^2/4$ terms only yields a lower bound on the approximation error of order
$2^{-\mu n^2/4}$, against a required bound of order $2^{-3n/2}$. The shortfall is a factor of
$\approx 0.42n$ \emph{in the exponent itself} — not a constant-factor gap that a sharper $\mu$
could close, but a gap that grows without bound as $n\to\infty$ for any fixed $\mu$. Zudilin's
bound $\mu(E)\le 2.42343562\ldots$ (Math.\ Notes 72 (2002) 858--862; erratum in Acta Arith.\ 111
(2004) 153--164 revises this to $\mu(E)\le 2.46497868\ldots$) is the best known value, but no future
improvement in $\mu$ can close this gap, because $\mu\ge 2$ always (a universal lower bound for any
irrational number), while the exponent shortfall grows like $n$ regardless of which fixed $\mu$ is
plugged in.

\textbf{Scope.} Closes: every attempt to derive the wall from a single-number irrationality-measure
statement about $C$ (or about $E$), present or future, as long as the measure is a fixed exponent
independent of $n$. Does not touch: simultaneous or multi-point Diophantine-approximation results
that could in principle scale with the number of terms (no such result is known to exist for this
constant, but the exponent-budget argument does not rule out the \emph{possibility} of one — only
the single-exponent route).
\coord{diophantine-approximation}

\textbf{Object or representation.} About the \emph{object}: this is a genuine value-theoretic
inequality (an exponent-budget computation), not an artifact of how the skip-set is coordinatized.

\textbf{Salvage.} Flagged in the source corpus as the single strongest reusable negative lemma in
the whole programme: any future proposal to attack a Mersenne/Erdős–Borwein-type constant via an
irrationality-measure bound should be checked against this exponent-budget argument \emph{first},
before any other work is invested.
\ev{Math}
\end{obs}

\begin{prop}[2-adic fixed-precision valuation-unit no-go — Lean-verified, problem-agnostic]
\label{prop:2adic-nogo}
\textbf{Route as conceived.} A carry orbit's local 2-adic signature — its valuation and odd unit
part at some fixed precision window — looks like a natural finite object to search for a
contradiction: if every carry orbit compatible with a given local valuation-unit word were forced
outside a safe centred interval, that would be a genuine local obstruction usable at every scale.

\textbf{Exact mechanism that closed it.} For \emph{any} finite word of odd-unit-part valuation
symbols at fixed precision $u>0$, and \emph{any} starting carry state $e$, there exists a
compatible carry orbit realizing that exact word with every intermediate state staying centred in
its dyadic interval, $|e'|\le \mathrm{vuRadius}\,u\,\sigma$. Bounded local valuation-unit data at
fixed precision can never exclude \emph{all} finite centred carry completions — the completion
always exists, constructively.
\[
  \exists\,\text{states},\; \mathrm{VUOrbit}\,u\,e\,\text{symbols}\,\text{states} \;\land\;
  (\forall\,\text{symbol},\, |\text{state}| \le \text{radius}).
\]
\lean{Erdos249257.TotientTailPeriodKiller.fixedPrecisionTropicalNoGo}{TropicalCurvatureCarry.lean:137}
built from
\lean{Erdos249257.TotientTailPeriodKiller.vu\_step\_has\_centred\_completion}{TropicalCurvatureCarry.lean:69}
and
\lean{Erdos249257.TotientTailPeriodKiller.vu\_word\_has\_prefix\_locked\_completion}{TropicalCurvatureCarry.lean:115}.

\textbf{Scope.} Closes: any proof strategy that tries to derive a contradiction purely from "the
local valuation-unit signature at fixed precision $u$ is incompatible with $X$" — such signatures
are \emph{always} realizable by some carry orbit, for every $u$ and every starting state. Does not
touch: arguments using \emph{growing} precision (an unbounded window, not a fixed one), or
arguments that couple the local signature to extra arithmetic (e.g.\ a global divisibility or
CRT constraint) rather than using the local signature alone.
\coord{2-adic-valuation-unit}

\textbf{Object or representation.} About a \emph{representation}: the theorem is entirely about
what a bounded local 2-adic window of a carry orbit can and cannot pin down; it says nothing about
the target value $C$ or the achievement set itself. It is a statement that a particular family of
finite coordinates (fixed-precision valuation-unit words) is too coarse to see the obstruction, not
that no obstruction exists.

\textbf{Salvage.} The theorem is stated and proved for \emph{any} starting carry state and any odd
unit-part word — zero arithmetic content beyond 2-adic bookkeeping. It is directly reusable as a
hard stop against any future fixed-window 2-adic attack on \emph{either} \#249 or \#257: a proof
needs growing precision or extra coupling, never a fixed local signature alone.
\ev{Lean}\scale{n/a}
\end{prop}

Read together, Observations~\ref{obs:prime-doubling}--\ref{obs:irrationality-measure} and
Proposition~\ref{prop:2adic-nogo} close every attack
family identified against the wall as of this writing that does not go through the digit-level
carry machinery itself. What remains open (Part~III, and \S\ref{sec:branch-exclusions} below) is
exactly the digit-carry route — this is not a residual gap after eliminating easier options; it is
the one route none of these four mechanisms can even in principle reach, because each closure
above excludes a \emph{different} kind of coarse-grained argument (density, tautological parity,
single-exponent approximation, fixed-window valuation) while leaving the exact digit recursion
itself untouched.

\subsection{Branch exclusions at the final producer: the exceptional dyadic cells}
\label{sec:branch-exclusions}

The two-sided dyadic invariant $\min(\mathrm{rem}(s), \mathrm{overshoot}(s)) \le 2^s$ (all $s\ge 5$) is the
sharpest general-purpose control on the seam orbit's scale in the whole corpus. It propagates by
induction from a single local socket, \textsf{SeamTwoSidedDyadicCellEscape}, which at every row
excludes a small number of \emph{exceptional cells} on the middle branch and imposes one further
inequality on the right branch:
\[
  \mathsf{SeamTwoSidedDyadicCellEscape} :=
  \forall\, s\ge 5,\;
  \Big(\lnot\mathrm{carries} \to \mathrm{middle} \to C_s \ne -3 \land C_s \ne -2 \land C_s \ne -1\Big)
  \;\land\;
  \Big(\lnot\mathrm{carries} \to \mathrm{right} \to \mathrm{overshoot}\le 2^s \to \mathrm{charge}\le 2^{s+2}\Big),
\]
where we write $C_s := 4\cdot\mathrm{rem}(s) - \mathrm{belowPulse}(s) - 4$ for the integer coordinate that
appears on the middle branch of the seam recursion (the same quantity that, along an all-right run,
is the affine excess $\mathrm{rem}(s)-2^s$ one step later — Part I, RC-2).
\lean{Erdos249257.SeamTwoSidedDyadicCellEscape}{HalfCylinderMiddleCarryLowerBound.lean:4374}

\begin{prop}[The upper (carries) successor needs no exceptional-cell exclusion]
\label{prop:upper-unconditional}
On the carries branch of the step (i.e.\ $(\mathrm{seamAdjacentCut}\ s\ hs).\mathrm{successorCarries}$
holds), the two-sided invariant propagates \emph{unconditionally}, with zero exceptional cells and
no hypothesis beyond \texttt{hcarry} itself:
\[
  \mathrm{successorCarries} \implies \mathrm{rem}(s+1) \le 2^{s+1}.
\]
\lean{Erdos249257.seamUpperBranch\_nextRemainder\_le\_pow}{HalfCylinderMiddleCarryLowerBound.lean:3924}
(from
\lean{Erdos249257.seamUpperBranch\_remainder\_add\_resetCharge\_eq}{HalfCylinderMiddleCarryLowerBound.lean:3542},
closed by \texttt{omega}). In this precise sense — no cell exclusion is ever needed on this branch
of the step — the upper successor contributes nothing to the exceptional-cell inventory: it is
"impossible" for the upper branch to be a source of an exceptional cell in the induction step,
because the branch is handled outright.
\coord{seam-branch-classification}\ \textbf{Object or representation:} about a
\emph{representation} (the induction-step case split), not the orbit itself.
\ev{Lean}\scale{n/a}
\end{prop}

\begin{thm}[$C_D=-3$ is impossible at the final middle producer, $D\ge 13$]
\label{thm:cd-neg3-impossible}
Consider a hypothetical \emph{final middle producer}: a middle transition at row $D\ge 13$
($\lnot\mathrm{carries}$, and the middle-branch inequality
$4\mathrm{rem}(D)+\mathrm{gap}-\mathrm{belowPulse} < \mathrm{terminalWeight}$ holds at $D$), followed
by an all-right tail forever after ($\forall s\ge D{+}1$, $\mathrm{seamGreedyWord}(s+1) =
\mathrm{seamGreedyWord}(s).\mathrm{extend}\ \mathrm{true}$). Then, unconditionally,
\[
  \mathrm{belowPulse}(D) + 2 \;\le\; 4\cdot\mathrm{rem}(D),
\]
i.e.\ $C_D \ge -2$ — the cell $C_D=-3$ (and every more negative value) is excluded.
\lean{Erdos249257.middleThenAllRight\_landingExcess\_two\_le}{HalfCylinderMiddleCarryLowerBound.lean:2319}

\textbf{Mechanism.} The proof chains three facts: (i) an all-right tail after $D$ forces the
terminal-augmented finite prefix at $D{+}1$ strictly above $1/2$
(\lean{Erdos249257.half\_lt\_upper\_competitor\_of\_eventually\_right}{HalfCylinderFatalGapRightTail.lean:627}),
hence the producer carry is strictly below its complete future incidence tail
(\lean{Erdos249257.middleProducer\_allRight\_forces\_carry\_lt\_tail}{HalfCylinderMiddleCarryLowerBound.lean:1814});
(ii) that carry/tail inequality transports to the exact rational model as
$\mathrm{seamGreedyFloorZ}(D) < \mathrm{seamTakeThreshold}(D)$
(\lean{Erdos249257.middleProducer\_allRight\_forces\_floorZ\_lt\_takeThreshold}{HalfCylinderMiddleCarryLowerBound.lean:1869});
(iii) combined with the pulse-absorption bound
$\mathrm{belowPulse}(D)\le 4\cdot\mathrm{seamWordFloorError}(D)$ and the exact floor identity
$\mathrm{seamGreedyFloorZ}(D)=\mathrm{rem}(D)-\mathrm{seamWordFloorError}(D)$, plus the unconditional strip
bound $1/3 < 4^D\cdot\mathrm{mersenneTail}(D)-2^D$, the strict inequality forces
$\mathrm{belowPulse}(D)+1 < 4\mathrm{rem}(D)$ in $\mathbb{Q}$, hence $+2\le$ in $\mathbb{N}$.

\textbf{Scope.} This closes $C_D=-3$ \emph{only inside the final-middle-producer-then-all-right
scenario} at $D\ge 13$ — not at every late middle row unconditionally. It is a strictly narrower
claim than \textsf{SeamTwoSidedDyadicCellEscape}'s general induction step, which still needs all
three cells excluded at \emph{every} late middle row (not just a hypothetical final one) to
propagate the two-sided invariant everywhere. The two statements share the same coordinate $C_D$
but differ in quantifier scope: Theorem~\ref{thm:cd-neg3-impossible} is a $\exists$-scenario
closure (one specific final producer under an all-right hypothesis), while the general socket is a
$\forall$-row closure. Do not conflate the two: closing $C_D=-3$ here does not close it in the
general induction step.
\coord{final-producer/landing-excess}\ \textbf{Object or representation:} about the \emph{object}
under the stated hypothesis — a genuine arithmetic consequence of the all-right-tail assumption
via the fatal-gap orbit, not a coordinate artifact. \ev{Lean}\scale{n/a}
\end{thm}

\begin{cor}[Only $C_D\in\{-2,-1\}$ remain among the exceptional dyadic cells]
\label{cor:cd-remaining}
Combining Proposition~\ref{prop:upper-unconditional} (upper branch: zero exceptional cells, ever)
with Theorem~\ref{thm:cd-neg3-impossible} ($C_D=-3$ closed, in the final-producer sub-case): of the
three cells $\{-3,-2,-1\}$ that \textsf{SeamTwoSidedDyadicCellEscape} must in general exclude on the
middle branch, only $C_D\in\{-2,-1\}$ remain genuinely open. No Lean theorem in the tree at the time
of writing excludes $C_D=-2$ or $C_D=-1$, in either the final-producer scenario or the general
induction step.
\coord{final-producer/landing-excess}\ \ev{Lean}
\end{cor}

\begin{rem}[The tail-dominance criterion — stated exactly, and not proved]
\label{rem:tail-dominance-open}
Write $\Theta_D$ for the real-valued producer-carry-to-tail ratio quantity that Theorem
\ref{thm:cd-neg3-impossible}'s mechanism step (i) actually compares against $C_D$: concretely,
$\Theta_D$ is the scaled complete-tail budget
$\mathrm{binaryCoeffTail}(\mathrm{supportCoeff}(\text{terminal-augmented prefix}))(2D+2)$ that
appears on the right-hand side of
\lean{Erdos249257.middleProducer\_allRight\_forces\_carry\_lt\_tail}{HalfCylinderMiddleCarryLowerBound.lean:1814},
expressed in the same integer units as $C_D$. The natural closing statement suggested by chaining
\textsf{SeamMiddleProducerCardEscape} / \textsf{SeamMiddleProducerRowEscape}
(\lean{Erdos249257.SeamMiddleProducerCardEscape}{HalfCylinderMiddleCarryLowerBound.lean:279},
\lean{Erdos249257.SeamMiddleProducerRowEscape}{HalfCylinderMiddleCarryLowerBound.lean:1769}) with
Corollary~\ref{cor:cd-remaining} is the \textbf{tail-dominance criterion}:
\[
  \textbf{(TD)}\qquad \forall\, D\ge 13,\ \big(\text{late middle row},\ C_D\ne -3\big) \implies
  \Theta_D < C_D.
\]
If (TD) held at every late middle row, it would supply exactly the missing half of
\textsf{SeamTwoSidedDyadicCellEscape}'s remaining two cells and close the whole
$\sqrt{}$-escape wall via
\lean{Erdos249257.half\_mem\_mersenneAchievementSet\_of\_middleProducerCardEscape}{HalfCylinderMiddleCarryLowerBound.lean:2364}.
\textbf{(TD) is NOT proved anywhere in the tree.} It does not appear as a theorem, a certified
computation, or even a numerically-checked conjecture in the files read for this survey — it is
this section's own naming of the gap the corollary leaves open, stated so that a future attempt has
an exact target rather than a vague "close the remaining cells" instruction. Anyone attempting it
should start from \textsf{SeamMiddleProducerRowEscape} (the weaker, row-scale sufficient condition
$s\le\mathrm{rem}(s)$), not from the sharper card-escape socket.
\ev{Open}\scale{cofinal}
\end{rem}

\subsection{No-go, countermodel, and rigidity modules that bound the wall from outside}

The digit-carry route is not merely unclimbed; a family of general-purpose no-go and countermodel
results independently constrains \emph{what kind of proof} can ever close it. These are not about
the wall's numerics — they are about which proof architectures are structurally impossible, for any
input.

\begin{prop}[Finite-state no-go: no bounded carry-state summary can recover history]
\label{prop:finite-state-nogo}
For a "balanced pulse" family at location $m$ (radius $r=(m{+}1)/2$, moving mass between positions
$m$ and $m{+}1$ without changing the series value), if a predecessor state is constant across the
whole family, then no function $\mathrm{decode}:\mathrm{State}\to\mathbb N$ can recover the
parameter $r$ from $\mathrm{state}(r)$ for every $r$ — fan-out is unbounded ($\ge\lfloor m/2\rfloor+2$
at $m=2k$). More strongly, any finite $\mathrm{Fintype}\ \mathrm{State}$ needs
$\mathrm{card}(\mathrm{State})\ge\mathrm{radius}+1$, unbounded in $m$.
\lean{Erdos249257.balancedPulse\_no\_autonomous\_decoder}{GenericTailOrbitRigidity.lean:247}
(fan-out lower bound at
\lean{Erdos249257.balancedPulse\_fanout\_unbounded}{GenericTailOrbitRigidity.lean:237}).

\textbf{Scope.} Rules out, for \emph{either} \#249 or \#257, any proof strategy that tries to define
a bounded/autonomous "carry state" summarizing pre-$m$ history and use it to exactly determine the
post-$m$ tail — completely general, no dependence on whether the coefficients are $\varphi$ or a
Möbius-support indicator.
\coord{binary-digit, generic}\ \textbf{Object or representation:} about a \emph{representation}
class (any finite-automaton encoding of history) — it says nothing about $C$ or $\mathcal A$
directly, only that this whole family of encodings is too weak.
\ev{Lean}\scale{n/a}
\end{prop}

\begin{prop}[Möbius-support countermodel: the natural negative-sign candidate overshoots]
\label{prop:mobius-nogo}
The signed Lambert identity $\sum_{d\ge1}\mu(d)/(2^d-1) = 1/2$ is exact. Writing
$N:=\{d:\mu(d)=-1\}$: $\sum_{d\in N}1/(2^d-1) = 1/2 + \sum_{d\in P}1/(2^d-1)$ where $P:=\{d\ge2:
\mu(d)=1\}$, and quantitatively $1/2 + 1/63 \le \sum_{d\in N}1/(2^d-1)$ (using the first positive
tail term $d=6$, $\mu(6)=1$) — the negative-Möbius support strictly \emph{overshoots} $1/2$ by at
least $1/63$.
\lean{MobiusSignSupportNoGo.half\_lt\_tsum\_negativeMobius}{MobiusSignSupportNoGo.lean:164}
(exact decomposition at
\lean{MobiusSignSupportNoGo.tsum\_negativeMobius\_eq\_half\_add\_positiveMobiusTail}{MobiusSignSupportNoGo.lean:111}).

\textbf{Scope.} Rules out exactly one natural candidate infinite Boolean support (the negative-Möbius
set) as a witness for $1/2\in\mathcal A$. A route-sufficiency no-go only — it says nothing about
whether \emph{some other} infinite Boolean support sums to $1/2$.
\coord{mobius-mersenne}\ \textbf{Object or representation:} about the \emph{object} — a genuine
value inequality for one specific candidate set, not a coordinate artifact. Margin $1/63$ is the
concrete number any repair attempt (adding/removing finitely many elements) must close or exceed.
\ev{Lean}\scale{fixed}
\end{prop}

\begin{prop}[No finite support ever hits $1/2$; finite certificates only ever certify death]
\label{prop:finite-boolSupport-and-onesided}
No finite positive-index Boolean support has value exactly $1/2$: the reduced denominator of any
finite Mersenne subset-sum is provably \textbf{odd} (each $2^n-1$ is odd), while $1/2$ needs an even
denominator.
\lean{finite\_boolSupport\_ne\_half}{HalfCarryReachability.lean:589}. Separately,
$\mathsf{CertifiedGreedyMersenneDeath}$ is a decidable, finite-depth certificate proving a real $x$
is \emph{not} in $\mathcal A$ (e.g.\ $3/4$ is machine-certified dead at level 1, lookahead 0:
\lean{three\_fourths\_certifiedGreedyMersenneDeath}{GreedyAchievementSet.lean:1742}), and this
one-sidedness is structural: \textbf{absence} of a found death certificate, or survival through any
finite depth, proves \emph{nothing} about membership.
\lean{Erdos249257.CertifiedGreedyMersenneDeath}{GreedyAchievementSet.lean:1720}.

\textbf{Scope.} The odd-denominator fact rules out every finite support as a $1/2$-witness,
unconditionally (this is exactly B4/C3 of Part~I/III's cut-locator machinery, and it is why any
$1/2$-achieving support, if one exists, must be genuinely infinite). The one-sidedness caveat rules
out treating a finite certified-kill search's silence as evidence of membership or rationality, for
\emph{any} member of the whole certified-kill family (this closure family, \#249's
\textsf{TotientTailPeriodKiller}, \textsf{LcmConeFlatness}, etc.) — a structural, not
problem-specific, warning that recurs across the entire certificate-based proof programme.
\coord{parity-denominator / finite-search}\ \textbf{Object or representation:} the odd-denominator
fact is about the \emph{object} (value inequality via denominator parity); the one-sidedness caveat
is about the \emph{representation} (what a finite search, as a proof method, can and cannot show).
\ev{Lean}\scale{fixed}
\end{prop}

\begin{prop}[Carry-orbit closure: \#249's rigidity engine, occupied ground for \#257]
\label{prop:carry-survivor-extinction}
For $\sum_n\varphi(n)/2^n$: a "survivor kill" certificate (every integer candidate in a bounded box
provably escapes a shrinking strip within $K$ steps, decidable) proves the tail difference
$\mathrm{totientTail}(N{+}h)-\mathrm{totientTail}(N)$ is not an integer. Rationality would force
integrality along an entire period ray, so one dead multiple $m\cdot h_0$ kills the whole primitive
period $h_0$
(\lean{Erdos257PeriodNoncollapse.irrational\_totient\_series\_of\_multiple\_survivor\_supply}{CarrySurvivorExtinction.lean:491}),
and this collapses onto the single $\mathbb N$-indexed family $\mathrm{periodLcm}(t)=\mathrm{lcm}(1,\ldots,t)$
(\lean{Erdos257PeriodNoncollapse.irrational\_totient\_series\_of\_lcm\_survivor\_supply}{CarrySurvivorExtinction.lean:540}).
Concrete unconditional deposit: every period $h\le16$ is machine-checked dead at $(N,L)=(14,9)$
(\lean{Erdos257PeriodNoncollapse.totient\_series\_ne\_rat\_of\_den\_dvd\_upto\_sixteen}{CarrySurvivorExtinction.lean:587}).

\textbf{Scope — a naming correction on record.} \textsf{CarrySurvivorExtinction.lean},
\textsf{AdjacentCarryTube.lean}, \textsf{AdjacentPhaseSeparation.lean}, and
\textsf{TotientCarryKernelRigidity.lean} are, despite names suggestive of \#257 support-rigidity,
\textbf{\#249 files} (namespace \texttt{Erdos249257.TotientTailPeriodKiller}, statements entirely
about $\sum\varphi(n)/2^n$). They are recorded in this closed-routes ledger not because they close
anything about \#257 directly, but because (i) their engine is the shared tempered-orbit machinery
also used by \#257's own rigidity results (Part~III's F1/T7), and (ii) a \#257-support reader must
know this ground is \emph{occupied by \#249}, not available — a route attempting to reuse these
exact declarations for a \#257 support-side argument would be reusing the engine, not the theorem.
\coord{integer carry-orbit, totient-specific instantiation}\ \textbf{Object or representation:}
about the \emph{object} ($\sum\varphi(n)/2^n$'s actual value), for \#249 — irrelevant by content, not
by mechanism, to \#257.
\ev{Lean}\scale{cofinal}
\end{prop}

\subsection{The scalar-localisation height obstruction}

\begin{lem}[Denominator complement survives scaling]
\label{lem:scalar-localization}
For $x:\mathbb Q$, $c:\mathbb Z$, $H:\mathbb N$: if $H\mid x.\mathrm{den}$ and
$(c\cdot x).\mathrm{den}\mid H$ — i.e.\ multiplying by the integer $c$ shrinks the displayed
denominator down \emph{into} $H$ — then the \textbf{complementary} denominator factor
$x.\mathrm{den}/H$ divides $c$:
\[
  H\mid x.\mathrm{den} \ \land\ (c\cdot x).\mathrm{den}\mid H \implies x.\mathrm{den}/H \mid |c|.
\]
\lean{Erdos249257.AdelicHeightObstruction.scalarLocalization\_complement\_dvd}{AdelicHeightObstruction.lean:23}.
Equality form (height conservation): $\exists\,t:\mathbb Z,\ H\cdot c\cdot x = t\cdot x.\mathrm{num}$
— the omitted denominator is transferred exactly to the coefficient $t$, never erased.
\lean{Erdos249257.AdelicHeightObstruction.scalarLocalization\_integer\_eq\_mul\_num}{AdelicHeightObstruction.lean:56}.

\textbf{Scope.} Fully generic $\mathbb Q$-arithmetic — no reference to either Erdős problem. It is
the reusable "clearing a denominator by multiplying by a small integer $c$ cannot discard the
complementary denominator factor; it only moves that factor into $c$" lemma. It gives a hard floor
on how small $c$ can be if it is meant to kill a target denominator factor $x.\mathrm{den}/H$: any
argument on either problem that hopes to "clear a denominator by a bounded multiplier" is bounded
below by this lemma.
\coord{rational-denominator/height}\ \textbf{Object or representation:} about a
\emph{representation} — pure denominator bookkeeping for $\mathbb Q$, with zero problem-specific
content; the upstream primitive
(\lean{Erdos249257.RationalDenominatorSurvival.divisor\_dvd\_divInt\_den}{RationalDenominatorSurvival.lean:17})
is the domain-neutral extraction this lemma builds on.
\ev{Lean}\scale{n/a}
\end{lem}

\begin{cor}[Mersenne specialisation]
\label{cor:mersenne-height}
For positive $x:\mathbb Q$: if $2^r\mid x.\mathrm{num}.\mathrm{natAbs}$ and $x < 2/(2^n-1)$, then
$2^r\cdot(2^n-1) < 2\cdot x.\mathrm{den}$ — a numerator $2$-power lower bound plus a Mersenne-scale
upper bound on $x$ together force a denominator lower bound, \emph{without} introducing a global
prefix LCM.
\lean{Erdos249257.AdelicHeightObstruction.positiveRat\_mersenne\_height}{AdelicHeightObstruction.lean:103}
(generic form
\lean{Erdos249257.AdelicHeightObstruction.positiveRat\_numDivisor\_mul\_lt\_two\_mul\_den}{AdelicHeightObstruction.lean:77}).
A ready-made denominator-lower-bound tool for "synchronized tail" arguments on \#257's Mersenne
sums; composes directly with Lemma~\ref{lem:scalar-localization} for a full numerator/denominator
height-transport toolkit.
\coord{rational-denominator/height, Mersenne instance}\ \ev{Lean}\scale{n/a}
\end{cor}

\subsection{The final-skip band formula, and what it does not show}

\begin{defn}[Upper-reset dyadic band escape]
\label{defn:band-escape}
\[
  \mathsf{SeamUpperResetDyadicBandEscape} :=
  \forall\, d\ge 13,\ \mathrm{successorCarries}(d) \to \forall\, j\le d,\quad
  2^{d-j+1} < 4\cdot\mathrm{overshoot}(d)+\mathrm{abovePulse}(d)
  \ \lor\
  4\cdot\mathrm{overshoot}(d)+\mathrm{abovePulse}(d)+2(d+j) \le 2^{d-j+1}.
\]
\lean{Erdos249257.SeamUpperResetDyadicBandEscape}{HalfCylinderMiddleCarryLowerBound.lean:4566}.
In words: at every upper-reset row $d\ge13$, the reset charge must avoid a linear-width forbidden
band immediately below every dyadic threshold $2^{d-j+1}$, for every $j\le d$ simultaneously.
Granted this socket,
\lean{Erdos249257.half\_mem\_mersenneAchievementSet\_of\_upperResetDyadicBandEscape}{HalfCylinderMiddleCarryLowerBound.lean:4790}
gives $1/2\in\mathcal A$ outright.
\end{defn}

\begin{prop}[Quantifier collapse: one critical index suffices, not $d{+}1$]
\label{prop:critical-band-index}
The abstract, purely combinatorial statement
$\mathsf{DyadicBandEscape}(d,E) \iff \exists j,\ \mathsf{CriticalDyadicBandIndex}(d,E,j)\land
E+2(d+j)\le 2^{d-j+1}$ collapses the $\forall j\in[0,d]$ band-avoidance condition (formally $d{+}1$
separate inequalities) to checking exactly \emph{one} nearest-boundary index $j$. Specialised to the
concrete seam reset charge, $\mathsf{SeamUpperResetCriticalBandEscape}$ is proved logically
\emph{equivalent} to Definition~\ref{defn:band-escape}'s socket.
\lean{Erdos249257.dyadicBandEscape\_iff\_exists\_critical}{HalfUpperResetCriticalBand.lean:108}
(seam specialisation
\lean{Erdos249257.seamUpperResetCriticalBandEscape\_iff}{HalfUpperResetCriticalBand.lean:883}).
Zero Mersenne/seam content in the core lemma — pure $(d,E,j)$ arithmetic over powers of 2.
\coord{dyadic-boundary, generic}\ \ev{Lean}\scale{n/a}
\end{prop}

\begin{obs}[Rows 13 through 31 are unconditionally certified]
\label{cert:band-through-31}
$\mathsf{seamUpperResetDyadicBandEscape\_through\_thirty}$ verifies Definition~\ref{defn:band-escape}
by kernel computation of $\mathrm{rem}(s)$ for every row $13\le s\le31$ — e.g.\ row 14 gives $\mathrm{rem}(14)=392$,
row 31 gives $\mathrm{rem}(31)=4187487147$ — with \textbf{zero} rows failing the band-avoidance test.
\lean{Erdos249257.seamUpperResetDyadicBandEscape\_through\_thirty}{HalfCylinderUpperResetBandCertificates.lean:78}.
\ev{Cert}\scale{bounded}
\end{obs}

\begin{rem}[What this does \emph{not} show — read the quantifiers exactly]
\label{rem:band-caveat}
Definition~\ref{defn:band-escape} is a $\forall d\ge13$ statement; Certificate
\ref{cert:band-through-31} is a finite, kernel-checked verification for $13\le d\le31$. \textbf{The
finite certificate does not establish the universal socket.} Nothing in the tree proves
$\mathsf{SeamUpperResetDyadicBandEscape}$ for all $d\ge13$ — this is exactly the open producer named
in Part~I (SE-9) and Part~III (RC-4/D1's \textsf{LargestSkipLateStepSocket}), stated here in its own
coordinate. The certified range through row 31 is real, certified evidence \emph{at that finite
range}, and it is a genuinely different fact from "the actual orbit avoids the unsafe band
cofinally" or "at every row." The socket remains an unproved sufficient condition; the finite
computation is not a proof of the cofinal or universal claim, and must never be quoted as one.
Concretely: the certificate says the band is avoided for $1264$ real resets through row $2500$ with
zero classification anomalies (Part~I, RC-7) and the sign law holds at all of them (RC-8) — this is
strong disconfirming-of-a-counterexample evidence, not a proof that the pattern continues past the
certified range.
\ev{Open}\scale{cofinal}
\end{rem}

